% =============================================================================
% REPLACEMENT FIGURE CAPTIONS
% =============================================================================

% Figure 1
\caption{\textbf{Empirical SNP-heritability point estimates, estimator-specific uncertainty, and cross-fold stability across analytical configurations.} Columns correspond to the ten empirical phenotypes, with binary traits shown on the primary liability scale and BMI on its native quantitative scale. Phenotype names are colour-coded consistently across panels. (A) Mean $h^2$ across available cross-validation folds. Values are retained without truncation; the colour scale is clipped at the 2nd--98th percentiles to prevent extreme values from obscuring the remaining configurations. (B) Median estimator-specific standard error across available folds. An $\times$ indicates that estimator-specific SE was not available in the parsed output. (C) Cross-fold SD, which quantifies sensitivity to the overlapping cross-validation training partitions and is interpreted as descriptive stability rather than inferential uncertainty. The SE and cross-fold-SD colour scales are capped at their 95th percentiles for visualisation; uncapped numerical values are reported in Supplementary Table~\ref{tab:configuration_uncertainty}.}

% Figure 2
\caption{\textbf{Distribution of empirical SNP-heritability estimates with corresponding uncertainty and cross-fold stability.} Each point represents one configuration--phenotype unit and is coloured by phenotype. Binary traits are shown on the liability scale and BMI on its native quantitative scale. (A) Distribution of mean $h^2$ values across analytical configurations; horizontal reference lines indicate $h^2=0$ and $h^2=1$. (B) Estimator-specific SE returned by the underlying software, shown on a logarithmic scale. LDpred2 points are absent from this panel because estimator-specific SE was not available in the parsed output. (C) Cross-fold SD across overlapping CV training partitions, shown on a symmetric logarithmic scale. Panels B and C represent different uncertainty concepts and should not be interpreted interchangeably.}

% Supplementary Figure S1
\caption{\textbf{Effect of target sample size on cross-fold stability in controlled simulations.} The same underlying simulated traits and analytical configurations were paired between $N=750$ and $N=5{,}000$. Points show method-level median cross-fold SD at each target sample size, and connecting lines show the direction of change. Binary phenotypes are evaluated on the liability scale and continuous phenotypes on their native quantitative scale. Benjamini--Hochberg-adjusted paired Wilcoxon $q$ values are shown for each method.}